\section{제안 방법}
\label{sec:method}

\subsection{전체 절차}
본 논문은 CoC 추론 주석을 안전/위험 레이블이나 제어 명령으로 직접 사용하지 않는다. 중심 절차는 연속 CoC를 행동 의도, 지속 의무, 해제 조건, 행동 순서와 삼값 증거로 변환하고 장면 내부 상태를 추적한다. 보조 절차는 같은 입력의 출처, 반복 정책, 잠정 제한 시간과 자차 움직임 반응을 별도 계약으로 검사한다. 두 절차는 다음 네 구성요소를 공유하되 결과를 서로 독립적으로 기록한다.

\begin{enumerate}
\item \textbf{증거 수집기(evidence collector)}는 \texttt{clip\_id} 안에서 타임스탬프와 CoC 텍스트를 수집하고, \texttt{quality\_pass} 및 자차 움직임을 함께 보존한다. 이 구성요소는 텍스트의 정답성이나 물리 위험을 판정하지 않고, 관측된 입력과 그 출처만 고정한다.
\item \textbf{계약 컴파일러(contract compiler)}는 수집된 텍스트에서 행동 의도, 지속 의무, 해제ㆍ충족 조건과 필수 행동 순서를 파싱한다. 중심 계약은 이 정보를 공통 중간표현으로 저장한다. 보조 반응 계약에는 잠정 제한 시간 $D_i$, 반복 정책(\preserve/\reset/\dedup)과 의도별 응답 술어 $R_{a_i}$를 추가한다. 이 규칙들은 데이터가 자동으로 발견한 안전 한계가 아니라 검증 전에 선택한 의미론이다.
\item \textbf{상태 관찰자(state observer)}는 시간순 계약 사건을 소비하여 의무의 생성--활성--충족 대기--해제 또는 위반 상태를 갱신한다. 보조 반응 계약에서는 응답 매칭, 저품질 사건의 검증 입력 수락 상태 덮어쓰기 금지와 장면 경계 종료도 검사한다. 이 상태는 단일 레이블의 분류 결과가 아니라 장면 내부 사건열의 프로토콜 상태이다.
\item \textbf{결과 분류기(result router)}는 검사 출력을 계약 통과, 의미론/임계값 의존 검토, 출처 부족, 인위적 오류 사례 결과로 구분한다. 계약 통과는 선택된 계약 아래 기록과 도출 응답이 일관적이라는 뜻이며, 검토 또는 출처 부족은 안전/위험 차량 레이블로 승격하지 않는다. 인위적 오류 사례 결과는 별도의 내부 타당성 확인 경로로만 보낸다.
\end{enumerate}

따라서 이 절차의 출력은 차량 제어가 아니라, 연속 CoC가 지정된 순서ㆍ상태 규칙을 만족하는지에 대한 판정과 위반 위치, 그리고 별도로 계산한 반응 시간ㆍ출처 검토 경로이다.

\begin{figure*}[t]
\centering
\resizebox{0.84\textwidth}{!}{%
\begin{tikzpicture}[
  box/.style={draw, rounded corners=2pt, align=center, minimum height=11mm,
    text width=2.55cm, font=\footnotesize},
  small/.style={draw, rounded corners=2pt, align=center, minimum height=8mm,
    text width=3.15cm, font=\scriptsize},
  obs/.style={box, fill=observed!10},
  der/.style={box, fill=derived!10},
  ass/.style={box, fill=assumed!10},
  regression/.style={box, fill=neutral!8},
  route/.style={box, fill=neutral!8},
  note/.style={small, fill=black!3},
  flow/.style={-{Latex[length=2mm]}, thick, draw=black!80},
  feedback/.style={-{Latex[length=1.8mm]}, dashed, thick, draw=black!55}
]

% Four-component evidence architecture.
\node[obs] (collector) at (0,0) {Evidence collector\\ordered CoC text\\timestamps, provenance};
\node[ass] (compiler) at (4.0,0) {Contract compiler\\obligation, release\\action order, unknown};
\node[der] (events) at (7.4,0) {Deterministic inputs\\common IR, XML\\queries};
\node[der] (observer) at (11.1,0) {State observer\\active obligation\\release, violation};
\node[route] (router) at (14.8,0) {Separated reports\\natural review\\auxiliary analyses};

\node[ass] (mut) at (7.4,-2.0) {Synthetic IR pairs\\one-field mutation\\enforced postcondition};
\node[der] (mutobserver) at (11.1,-2.0) {Event-local/stateful\\rule execution\\same compiled input};
\node[regression] (regression) at (14.8,-2.0) {Regression coverage\\not detection accuracy};
\node[note] (boundary) at (0,-2.0) {No environment search\\or control synthesis};

\node[route] (select) at (11.1,-4.0) {Selected natural sample\\blinded review\\separate consensus};
\node[note] (scope) at (7.4,-4.0)
  {Separate observed, derived, assumed, and mutated evidence.};

\draw[flow] (collector) -- (compiler);
\draw[flow] (compiler) -- (events);
\draw[flow] (events) -- (observer);
\draw[flow] (observer) -- (router);
\draw[feedback] (events) -- (mut);
\draw[flow] (mut) -- (mutobserver);
\draw[flow] (mutobserver) -- (regression);
\draw[flow] (router.east) -- ++(1.8,0) |- (select.east);
\draw[flow] (select) -- (scope);

% Feedback is routed outside the node area to avoid crossing boxes.
\coordinate (feedbackleft) at ([xshift=-6mm]boundary.west);
\draw[feedback] (select.south) -- ++(0,-0.6) coordinate (feedbackbottom)
  -- node[midway,below=1mm,font=\scriptsize,align=center]{Review feedback}
  (feedbackbottom -| feedbackleft) -- (feedbackleft) |- (collector.west);

\end{tikzpicture}%
}

\bifigcaption{증거 수집기, 계약 컴파일러, 상태 관찰자와 결과 분류기로 구성한 제안 방법. 선택한 자연 표본의 사람 검토와 사후조건을 강제한 합성 중간표현 회귀 시험은 서로 다른 증거 경로로 분리한다.}{Proposed method separating selected natural-sample review from postcondition-enforced synthetic-IR regression.}
\label{fig:method-pipeline}
\end{figure*}

그림~\ref{fig:method-pipeline}에서 증거 수집기는 관측값을 보존하고, 계약 컴파일러는 그 값을 상태 계약으로 변환하며, 상태 관찰자는 계약 상태를 검사하고, 결과 분류기는 판정의 사용 범위를 분리한다. 자연 표본의 검토 결과는 합성 변형을 만들거나 검사기 임계값을 맞추는 데 사용하지 않으며, 합성 회귀 시험 결과도 자연 표본의 정답으로 전용하지 않는다.

\subsection{UPPAAL 네트워크}
\uppaal 모델은 단일 오토마타가 아니라 실험 목적마다 다른 템플릿 네트워크이다. \texttt{scene-0001} 모델에서는 증거 수집기가 고정한 타임스탬프, 텍스트, 품질, 자차 움직임과 계약 컴파일러가 만든 사건 배열을 사용한다. \texttt{CoCSource}는 검증 입력 수락 시작, 반복, 검증 입력 비수락 사건을 타임스탬프 순서로, \texttt{EgoTrace}는 자차 움직임에서 도출된 차량 반응을 그 순서대로 재생한다. 이 원천 오토마타는 \emph{결정적 기록 사건열}을 한 번 재생할 뿐, 차량 환경을 탐색하지 않으며 안전 제어를 생성하거나 제어 정책을 합성하지 않는다. \texttt{ObligationManager}, \texttt{ReactionObserver}, \texttt{DecisionValidator}는 같은 사건열에서 의무 생명주기, 응답 매칭, 반복 시계 정책과 저품질 덮어쓰기를 검사한다. 분석 대상 집합과 체인 실험은 뒤의 표와 같이 별도 모델군이며, 하나의 통합 모델이 모든 실험에 실행된 것은 아니다.

\begin{figure*}[t]
\centering
\resizebox{0.78\textwidth}{!}{%
\begin{tikzpicture}[
  node distance=13mm and 15mm,
  box/.style={draw, rounded corners=2pt, align=center, minimum height=13mm,
    text width=3.25cm, font=\footnotesize},
  source/.style={box, fill=observed!12},
  monitor/.style={box, fill=derived!11},
  compiler/.style={box, fill=assumed!13},
  bus/.style={draw, rounded corners=2pt, fill=black!4, align=center,
    minimum width=11.7cm, minimum height=11mm, font=\footnotesize},
  flow/.style={-{Latex[length=2.3mm]}, line width=0.7pt, draw=black!75},
  event/.style={-{Latex[length=2.3mm]}, line width=0.85pt, draw=derived!80},
  audit/.style={-{Latex[length=2.3mm]}, line width=0.85pt, draw=assumed!85},
  label/.style={font=\scriptsize, align=center, fill=white, inner sep=1.5pt}
]

% Sources.
\node[source] (coc) at (0,3.4) {
  \textbf{CoCSource}\\
  replay: collector-fixed CoC\\
  intent-tagged decision events\\
  accepted (code: trusted) / repeat / untrusted
};

\node[source] (ego) at (7.7,3.4) {
  \textbf{EgoTrace}\\
  replay: collector-fixed ego motion\\
  selected $R_a$ applied\\
  derived response evidence
};

% Researcher-selected contract semantics, kept distinct from recorded evidence.
\node[compiler] (compiler) at (3.85,5.45) {
  \textbf{Contract compiler}\\
  selected intent, $D$, repeat, $R_a$\\
  contract semantics only
};

% Broadcast bus.
\node[bus] (bus) at (3.85,1.45) {
  \textbf{Deterministic replay bus}\\
  compiled decision events + derived ego response evidence\\
  \texttt{first!}, \texttt{repeat!}, \texttt{response!}, \texttt{untrusted!}
};

% Monitors.
\node[monitor] (obl) at (-1.6,-0.95) {
  \textbf{ObligationManager}\\
  \textbf{State observer}\\
  Idle $\rightarrow$ Active\\
  Active $\rightarrow$ Responded
};

\node[monitor] (react) at (3.85,-0.95) {
  \textbf{ReactionObserver}\\
  \textbf{State observer}\\
  clock $reaction$\\
  deadline $D$\\
  repeat policy
};

\node[monitor] (valid) at (9.3,-0.95) {
  \textbf{DecisionValidator}\\
  \textbf{State observer}\\
  quality gate\\
  overwrite guard
};

% Compiler selects contract semantics; sources replay the resulting evidence.
\draw[audit] (compiler.south west) -- node[label, left] {selected\\contract} (coc.north);
\draw[audit] (compiler.south east) -- node[label, right] {response\\predicate} (ego.north);

% Source to bus.
\draw[event] (coc.south) -- node[label, left] {compiled decision\\events} (bus.north west);
\draw[event] (ego.south) -- node[label, right] {derived ego\\response evidence} (bus.north east);

% Bus to monitors. Route from distinct bottom anchors to avoid visual crossing.
\draw[event] ([xshift=-4.2cm]bus.south) -- node[label, left] {\texttt{first?}\\\texttt{response?}} (obl.north);
\draw[event] (bus.south) -- node[label, right] {\texttt{first?}\\\texttt{repeat?}\\\texttt{response?}} (react.north);
\draw[event] ([xshift=4.2cm]bus.south) -- node[label, right] {\texttt{untrusted?}} (valid.north);

% Result box.
\node[draw, rounded corners=2pt, fill=neutral!8, align=center,
  minimum width=11.7cm, minimum height=10mm, font=\footnotesize] (queries)
  at (3.85,-3.05) {
  \textbf{Result router}: contract pass / review / provenance / mutation\\
  UPPAAL queries:
  \texttt{Active}, \texttt{Responded}, \texttt{DeadlineMiss},
  \texttt{clock\_reset}, \texttt{untrusted\_overwrite}, \texttt{deadlock}
};

\draw[flow] (obl.south) -- (queries.north west);
\draw[flow] (react.south) -- (queries.north);
\draw[flow] (valid.south) -- (queries.north east);

\end{tikzpicture}%
}

\bifigcaption{\texttt{scene-0001-normal}에 대응하는 \uppaal 네트워크 개념도. 증거 수집기와 계약 컴파일러의 결정적 사건열을 상태 관찰자가 소비하고, 결과 분류기가 질의 판정을 분리한다.}{Conceptual UPPAAL network corresponding to scene-0001-normal.}
\label{fig:uppaal-network}
\end{figure*}

그림~\ref{fig:uppaal-network}에서 위쪽 두 템플릿은 증거 수집기와 계약 컴파일러가 준비한 입력 흐름을 재생하는 원천이고, 아래쪽 세 템플릿은 상태 관찰자를 구성한다. 이 구분이 중요하다. \texttt{CoCSource}와 \texttt{EgoTrace}는 새로운 주행 행동을 탐색하지 않으며, 관찰자들은 이미 주어진 사건열 위에서 검증 항목 생성, 응답 매칭, 기한 실패, 출처 덮어쓰기와 경계 정책을 검사한다. 질의 결과의 의미는 결과 분류기가 계약 통과, 의미론/임계값 의존 검토, 출처 부족, 인위적 오류 사례로 분리한다.

그림~\ref{fig:uppaal-network}의 블록도는 표~\ref{tab:ta-skeleton}의 시간 오토마타 네트워크로 구현된다. 시간 단위는 XML 내부에서 밀리초 정수이며, 표에서는 \texttt{first}, \texttt{repeat}, \texttt{response}, \texttt{untrusted}를 각각 \texttt{first\_trigger}, \texttt{repeat\_trigger}, \texttt{response}, \texttt{untrusted\_event} 방송 채널의 약어로 사용한다. 원천 템플릿은 사건을 한 번 재생하고, 관찰자 템플릿들은 같은 방송 사건에 동기화되어 의무의 생성, 기한, 반복 의미론, 낮은 품질 사건의 덮어쓰기 금지를 독립적으로 검사한다.

\begin{table}[!htbp]
\centering
\bitablecaption{그림~\ref{fig:uppaal-network}에 대응하는 시간 오토마타 골격.}{Timed-automata skeleton corresponding to the UPPAAL network.}
\label{tab:ta-skeleton}
\scriptsize
\setlength{\tabcolsep}{2.5pt}
\renewcommand{\arraystretch}{0.90}
\begin{tabularx}{\columnwidth}{L{0.26\columnwidth}L{0.21\columnwidth}Y}
\toprule
Template & Locations and clocks & Main transitions and meaning \\
\midrule
\texttt{CoCSource} &
\texttt{C0--C4}, \texttt{episode} &
$0$s \texttt{first!}; $1.0$s, $2.3$s \texttt{repeat!}; $6.8$s \texttt{untrusted!}. Replays CoC decision events in timestamp order. \\
\midrule
\texttt{EgoTrace} &
\texttt{E0--E1}, \texttt{episode} &
$2.16$s \texttt{response!}. Replays weak response evidence derived from ego speed. \\
\midrule
\texttt{ObligationManager} &
\texttt{Idle}, \texttt{Active}, \texttt{Responded} &
Opens an obligation on \texttt{first?}; closes it on \texttt{response?}. Deadline state is tracked separately. \\
\midrule
\texttt{ReactionObserver} &
\texttt{Idle}, \texttt{Armed}, \texttt{Done}, \texttt{DeadlineMiss}; clock \texttt{reaction} &
Sets \texttt{reaction=0} on \texttt{first?}; the contract requires response before \texttt{reaction}$=D$, where the miss transition becomes enabled. \preserve does not restart \texttt{reaction} on repeats. \\
\midrule
\texttt{DecisionValidator} &
\texttt{Clean}; Boolean overwrite flag &
Keeps one location and sets the Boolean on \texttt{untrusted?} when the mutation enables an overwrite. \\
\bottomrule
\end{tabularx}
\end{table}

표~\ref{tab:ta-skeleton}에서는 그림~\ref{fig:uppaal-network}의 각 블록이 실제 모델에서 어떤 위치와 전환으로 구현되는지 정리한다. 예를 들어 \texttt{CoCSource}의 \texttt{first!} 전환은 검증 항목 생성을 촉발하고, \texttt{EgoTrace}의 \texttt{response!} 전환은 자차 속도에서 도출한 차량 반응을 방송한다. \texttt{ReactionObserver}는 이 두 사건 사이의 지역 시계 \texttt{reaction}을 검사하며, \preserve 기준 모델에서는 반복 사건을 보더라도 \texttt{reaction}을 재시작하지 않는다. 저장 집계에는 구현 시계가 기한과 같은 \texttt{reaction}$=D$인 사례가 없으며, 등호 시점의 전환 우선순위를 직접 고정하는 검사기 회귀 시험은 후속 과제로 남긴다.

\begin{table}[!htbp]
\centering
\bitablecaption{실험별 \uppaal 모델군. 중심 연속 CoC 관찰자는 시간 지연이 아니라 사건 순서와 의무 상태를 검사한다.}{UPPAAL model families used by the experiments; the central sequential-CoC observer checks event order and obligation state rather than elapsed-time bounds.}
\label{tab:model-families}
\scriptsize
\begin{tabularx}{\columnwidth}{L{0.25\columnwidth}L{0.30\columnwidth}Y}
\toprule
Experiment & Templates & Purpose \\
\midrule
Sequential CoC canonical fixtures & \texttt{ObserverTemplate} & One-template committed observer for ordered obligation/release semantics and Python-conformance checks \\
\midrule
\texttt{scene-0001} & \texttt{CoCSource}, \texttt{EgoTrace}, \texttt{ObligationManager}, \texttt{ReactionObserver}, \texttt{DecisionValidator} & Five-template timed lifecycle and mutation check \\
\midrule
94-scene batch & \texttt{EpisodeSource}, \texttt{LineageMonitor} & Two-template flag/query serialization regression \\
\midrule
13-chain smoke artifact & \texttt{LoopSource}, \texttt{Monitor} & Two-template implementation pipeline smoke test with a distinct parser/detector \\
\bottomrule
\end{tabularx}
\end{table}

이 골격에서 \texttt{CoCSource}와 \texttt{EgoTrace}는 환경을 새로 생성하는 반응형 제어기가 아니라, 관측 및 도출 증거 흐름을 결정적으로 재생하는 원천 오토마타이다. 반대로 \texttt{ObligationManager}, \texttt{ReactionObserver}, \texttt{DecisionValidator}는 같은 사건 흐름을 서로 다른 관점에서 감시하는 상태 관찰자 오토마타이다. 따라서 시간 오토마타의 검증 대상은 ``차량이 안전하게 주행했는가''가 아니라 ``검증 입력으로 수용된 CoC 주석 사건이 후보 의무를 열었는가'', ``반복 사건이 시계를 부당하게 초기화하지 않았는가'', ``차량 반응이 잠정 제한 시간 안에 매칭되었는가'', ``저품질 사건이 수용된 상태를 덮어쓰지 않았는가'', ``연속성 증거가 없는 경계에서 대기 의무가 닫혔는가''이다.

단일 결정적 재생만 실행한다면 동일한 기한 검사는 일반 스크립트로도 계산할 수 있다. 본 논문에서 모델 체킹을 사용하는 이유는 각 모델군의 상태 관찰자, 반복 의미론, 경계 정책 또는 변이 오라클을 선언하고, 위반 시 어떤 관찰자가 어떤 시각에 실패했는지 반례 궤적으로 얻기 위해서이다. 특히 \texttt{scene-0001} 모델은 같은 원천 사건을 의무 생명주기, 기한 관찰자와 출처 검증기가 동시에 소비하도록 만들기 때문에, 단순 수치 계산보다 ``어떤 정책 조합에서 어떤 프로토콜 속성이 깨졌는가''를 명세 수준에서 검토할 수 있다.

\subsection{scene-0001의 의미론 변환}
\texttt{scene-0001}의 원본 및 도출 정보는 표~\ref{tab:scene1-map}처럼 \uppaal 사건으로 변환된다. 표에서는 가독성을 위해 원본 타임스탬프를 장면 내부 초 단위로 표시하고, \uppaal 사건 시각은 최초 검증 입력 수락 시작 사건을 기준으로 한 상대 초 단위로 표시한다. 실제 XML 생성 시에는 이 값을 밀리초 정수로 양자화한다.

\begin{table}[!htbp]
\centering
\bitablecaption{\texttt{scene-0001} 증거를 \uppaal 사건으로 변환한 예. 시간 단위는 초이다.}{Mapping scene-0001 evidence to UPPAAL events; times are in seconds.}
\label{tab:scene1-map}
\footnotesize
\begin{tabularx}{\columnwidth}{L{0.30\columnwidth}L{0.25\columnwidth}Y}
\toprule
Source/derived evidence & \uppaal event & Meaning \\
\midrule
$t=2.500$s, verification-accepted \textsc{Decelerate} & \texttt{first!} at $0.000$s & Opens the first deceleration obligation \\
$t=3.500$s, repeated \textsc{Decelerate} & \texttt{repeat!} at $1.000$s & Observes a repeat of the same obligation \\
$t=4.660$s, speed-decrease onset & \texttt{response!} at $2.160$s & Weak response derived from ego trajectory \\
$t=4.800$s, repeated \textsc{Decelerate} & \texttt{repeat!} at $2.300$s & Repeated label after response evidence \\
$t=9.300$s, low-quality CoC & \texttt{untrusted!} at $6.800$s & Code-named event that must not overwrite verification-accepted state \\
\bottomrule
\end{tabularx}
\end{table}

표~\ref{tab:scene1-map}의 \texttt{first!}, \texttt{repeat!}, \texttt{untrusted!}는 각각 생성 XML의 \texttt{first\_trigger!}, \texttt{repeat\_trigger!}, \texttt{untrusted\_event!} 채널을 줄여 표기한 것이다.

생성되는 대표 질의는 다음과 같다.

\noindent\begin{minipage}{\columnwidth}
\begin{lstlisting}[caption={대표 \uppaal 질의 집합.},label={lst:queries}]
E<> Obligation.Active
E<> Obligation.Responded
A[] not Reaction.DeadlineMiss
A[] not clock_reset_violation
A[] not untrusted_overwrite
A[] not orphan_response
A[] not boundary_violation
A[] not deadlock
\end{lstlisting}
\end{minipage}

본 연구의 \uppaal 모델은 폐루프 반응형 제어기가 아니다. 기록된 CoC 사건 흐름과 자차 움직임 응답 흐름을 재생하고 관찰자 오토마타가 프로토콜 수준 속성을 검사한다. 따라서 \uppaal은 환경을 탐색하거나 안전 제어를 생성하지 않는다. 단일 결정적 재생 하나의 PASS만으로는 모델 체킹 의미가 크지 않다. 본 논문에서 \uppaal의 가치는 여러 관찰자, 반복 의미론, 변이 오라클, 다중 장면 경계 정책을 같은 정형 프레임워크 안에서 선언적으로 결합하고 반례 궤적을 얻는 데 있다.

\subsection{인위적 오류 사례 오라클}
인위적 오류 사례 오라클은 결과 분류기가 정상 계약 통과와 알려진 프로토콜 결함을 구분하는지 확인하는 \emph{검사기 회귀 시험} 장치이다. 기준 모델의 PASS가 질의 비활성화나 모델 구성 오류로 생긴 결과가 아님을 점검하기 위해, 정상 궤적과 함께 의도적으로 결함을 삽입한 변이 모델을 만들고 각 변이의 기대 실패 질의와 비교한다. 이 검사는 실제 오류 검출률, 실제 주행 결함의 발생 빈도 또는 미관측 배치 행동을 추정하지 않는다.

\begin{table}[!htbp]
\centering
\bitablecaption{변이 오라클과 기대 실패 속성.}{Mutation oracles and their expected failing properties.}
\label{tab:mutation-oracle}
\footnotesize
\begin{tabularx}{\columnwidth}{L{0.28\columnwidth}YY}
\toprule
Mutation & Injected fault & Expected failure \\
\midrule
\texttt{deadline\_2s} & Shrink deadline while keeping response & Deadline query \\
\texttt{no\_response} & Remove ego response & Response reachability, deadline \\
\texttt{clock\_reset} & Restart clock on repeat & Clock-reset query \\
\texttt{untrusted\_ovr.} & Low-quality event overwrites audit-accepted state & Provenance query \\
\texttt{orphan\_response} & Remove obligation trigger but keep response & Orphan-response query \\
\texttt{late\_stop} & Delay stop/hold stage & Stage deadline \\
\texttt{premature\_rec.} & Recover before hazard clearance & Clearance precondition \\
\bottomrule
\end{tabularx}
\end{table}

표~\ref{tab:mutation-oracle}의 \texttt{untrusted\_ovr.}와 \texttt{premature\_rec.}는 각각 \texttt{untrusted\_overwrite}와 \texttt{premature\_recovery} 변이를 축약한 것이다. 변이는 실제 데이터 결함이라고 주장하기 위한 것이 아니다. 검사기와 질의 집합이 알려진 프로토콜 결함을 구분할 수 있는지 확인하기 위한 합성 시험이며, 정상 궤적 지도 신호나 실제 오류 검출률의 근거로 사용해서는 안 된다.

\subsection{다중 장면 확장}
로컬 CoC-Nusc 단독 데이터에는 장면 간 전역 타임스탬프, 로그 식별자, 이전/다음 관계, 전역 자차 자세가 부족하다. 따라서 기본 정책은 모든 장면을 독립 에피소드로 닫고 장면 경계를 넘어 대기 의무를 이월하지 않는 것이다. 이 정책은 오류 차단 시험 실험으로 사용된다. 연속성 증거가 없는 경계에서 의무가 다음 장면까지 살아 있으면 변환 절차의 오류로 판정한다.

연속성이 입증된 경우에는 여러 장면을 하나의 장기 에피소드로 연결할 수 있다. 예비 실험에서는 nuScenes 호환 메타데이터 미러를 사용해 같은 로그의 인접 장면 쌍을 찾고, 간격이 1초 이하이며 로컬 추론/자차 움직임이 있는 장면들을 체인으로 묶었다. 이후 각 장면 내부 사건 타임스탬프에 전역 장면 오프셋을 더해 하나의 시간선을 만들고, 루프 인덱스 방식의 \uppaal 모델로 검증했다. 루프 인덱스 모델은 사건마다 전환을 모두 쓰는 대신 \texttt{idx}, \texttt{event\_time[]}, \texttt{event\_kind[]}, \texttt{event\_gap[]} 배열을 사용하므로, 장면 수가 늘어나도 오토마타 구조는 유지되고 데이터 배열만 커진다.

\subsection{연속 CoC 계약과 삼값 증거}
\label{sec:sequential-contract-method}

앞의 시간ㆍ출처 검증과 별도로, 같은 장면의 CoC 문장들이 만드는 요구 관계를 공통 계약 중간표현으로 컴파일한다. 사건 $e_i$는 요구 행동, 지속 의무, 해제 조건, 선행 행동과 출처 식별자를 가진다. 컴파일러는 ``정지한다'', ``통과할 때까지'', ``그 뒤 진행한다''와 같은 단계형 절을 하나의 최종 행동으로 축약하지 않고 순서대로 보존한다. 행동은 \textsc{Stop/Hold}, \textsc{Yield/Decelerate}, \textsc{Accelerate/Proceed}, \textsc{Maintain}의 추상 범주로 정규화한다.

해제와 객체 증거는 \textsc{True}, \textsc{False}, {\unknown}의 삼값으로 처리한다. 명시적 해제 증거가 참이면 의무를 종료하고, 거짓이면 의무를 유지한다. 문장이나 관측 자료가 해제 여부를 제공하지 않으면 의무 상태를 임의로 바꾸지 않고 {\unknown} 사유를 기록한다. 이 규칙은 증거 부족을 정상 사례로 계산하는 것을 막는다.

상태형 검사기는 장면별 활성 의무 집합을 유지하며 다음 네 모순을 검사한다. \texttt{HOLD\_GO\_CONFLICT}는 유지 의무가 활성인 동안 진행ㆍ가속이 시작된 경우, \texttt{PREMATURE\_RELEASE}는 해제 조건이 거짓인데 의무를 종료한 경우이다. \texttt{ORDER\_VIOLATION}은 필수 정지ㆍ감속 단계 전에 후속 행동을 시작한 경우이고, \texttt{STALE\_OBLIGATION}은 해제 조건 뒤에도 의무가 남아 진행성을 잃은 경우이다. 사건별 기준선은 같은 컴파일 결과를 사용하지만 이전 사건의 활성 의무를 다음 사건으로 넘기지 않는다.

동일한 사건 배열과 삼값 증거는 단일 \texttt{ObserverTemplate}의 \uppaal 관찰 모델로 변환된다. 이 템플릿은 committed 위치에서 배열 인덱스를 증가시키며 의무의 생성ㆍ유지ㆍ해제ㆍ위반 상태를 갱신한다. 위반 플래그와 최초 위반 사건 인덱스를 고정하여 \texttt{A[] !violation} 안전 질의와 위반 상태 도달성 질의를 검사한다. 이 중심 모델은 실제 타임스탬프 간 지연을 제약하지 않고 사건 순서만 재생한다. 따라서 앞의 반응 시간 모델과 달리 시간 기한을 검증하지 않으며, Python 판정이 어떤 활성 의무와 사건 순서를 거쳐 발생했는지 설명하는 역할만 한다.

40쌍의 회귀 시험은 자연 CoC 문장을 최소 수정한 표본이 아니라 \texttt{CONTROLLED\_SYNTHETIC\_IR} 출처의 결정적 합성 중간표현이다. 각 기준 사례에서 한 필드만 바꾸며, 생성기는 기준 사례가 상태형 검사에서 모순이 아니고 변형 사례가 지정한 모순 하나로 판정되는지를 사후조건으로 강제한다. 그러므로 40/40은 검사기가 자신의 선언된 규칙과 시험 fixture를 일관되게 처리했다는 커버리지 기록이며, 독립 평가셋의 검출 성능이나 통계적 일반화 근거가 아니다.
